\documentclass[11pt]{article}

\usepackage[T1]{fontenc}
\usepackage[utf8]{inputenc}
\usepackage{amsmath,amssymb,mathtools,mathrsfs}
\usepackage[margin=1in]{geometry}
\usepackage{microtype}
\usepackage{needspace}

\setlength{\parindent}{0pt}
\setlength{\parskip}{0.72em}

\begin{document}

\section*{How the Proof Works}

The proof begins with the way a smooth three-dimensional flow could fail. If its maximal time \(T_*\) were finite, Navier--Stokes would have to drive the velocity into a singularity:
\[
\partial_tu+(u\cdot\nabla)u+\nabla p=\nu\Delta u,
\qquad \nabla\cdot u=0,
\qquad u(\cdot,0)=u_0.
\]
The danger is concentration. Nonlinear transport can press variation into a smaller scale, while viscosity diffuses it. A vortex tube shows both effects acting on one piece of fluid. The surrounding strain lengthens the tube. Incompressibility prevents that lengthening from creating volume, so its cross-section contracts and the rotating core narrows.

The strain also acts on the vorticity \(\omega=\nabla\times u\). Curling the equation gives
\[
\partial_t\omega+(u\cdot\nabla)\omega
=(\omega\cdot\nabla)u+\nu\Delta\omega.
\]
Where the vortex is aligned with an amplifying direction of the strain,
\[
\omega\cdot((\omega\cdot\nabla)u)>0,
\]
the stretching part raises \(\lvert\omega\rvert\) as the core narrows. Since vorticity is one spatial derivative of velocity, the velocity changes across less distance and its gradients steepen. This is the possible route toward singularity: the nonlinear flow sharpens a structure at the very moment viscosity is trying to smooth it.

The question is which effect gains strength as the width shrinks. If a feature of width \(r\) and velocity scale \(U\) is compressed by a factor \(\lambda>1\), its new width is \(r/\lambda\) and its velocity scale is \(\lambda U\). The nonlinear and viscous terms then change as
\[
\frac{U^2}{r}
\longmapsto
\lambda^3\frac{U^2}{r},
\qquad
\nu\frac{U}{r^2}
\longmapsto
\lambda^3\nu\frac{U}{r^2}.
\]
They strengthen in exact proportion. This is the Navier--Stokes scaling. From a solution \((u,p)\), form
\[
u_\lambda(x,t)=\lambda u(\lambda x,\lambda^2t),
\qquad
p_\lambda(x,t)=\lambda^2p(\lambda x,\lambda^2t).
\]
Substitution multiplies every term in the equation by \(\lambda^3\). Compression does not make viscosity outrun the nonlinear term. Nor does it make the nonlinear term outrun viscosity. It preserves their contest on the smaller scale.

Kinetic energy does not see this concentration sharply enough. The compressed core occupies \(\lambda^{-3}\) of the volume while the square of its velocity contributes \(\lambda^2\), so
\[
\|u_\lambda(\cdot,t)\|_2^2
=\lambda^{-1}\|u(\cdot,\lambda^2t)\|_2^2.
\]
The core can grow narrower and its gradients can grow larger while its contribution to the energy falls. Across the homogeneous Sobolev scale,
\[
\|u_\lambda(\cdot,t)\|_{\dot H^s}
=\lambda^{s-\frac12}
\|u(\cdot,\lambda^2t)\|_{\dot H^s}.
\]
Concentration makes the norm smaller below \(s=\tfrac12\) and larger above it. At \(s=\tfrac12\), it does neither. The critical height is the scale at which narrowing the flow gives the measurement no automatic gain and no automatic loss. Write
\[
E_s(t):=\frac12\|\Lambda^su(t)\|_2^2,
\qquad
H(t):=E_{1/2}(t).
\]

The clock also contracts. A feature compressed by \(\lambda\) lives on a time scale shorter by \(\lambda^2\). If an evolving core were to repeat the contraction by a fixed factor \(0<\rho<1\), its successive scales would be
\[
r_j=\rho^jr_0,
\qquad
U_j=\rho^{-j}U_0,
\qquad
\tau_j=\rho^{2j}\tau_0,
\]
and
\[
\sum_{j=0}^{\infty}\tau_j
=\tau_0\sum_{j=0}^{\infty}\rho^{2j}
=\frac{\tau_0}{1-\rho^2}<\infty.
\]
Infinitely many narrower and faster stages fit inside a finite interval. This is the Zeno cascade. It is not a construction of a singular solution. It identifies what the proof must defeat: finite energy does not prevent concentration, and scaling gives viscosity no advantage over nonlinear sharpening as the scales collapse.

Suppose, for contradiction, that such a collapse ends the maximal smooth solution at \(T_*<\infty\).

In order for a singularity to form, velocity has to blow up. In order for that to happen, its acceleration would also have to blow up, which means its jerk has to blow up, and so on, and so on:
\[
u,\qquad \partial_tu,\qquad \partial_t^2u,\qquad \ldots
\]
What you would expect to see in this pattern, if a singularity were to form, is less spatial regularity as you go up the time derivatives, not more. The equation couples the two directions in the opposite way.

The first three rows are
\[
\partial_tu
=\nu\Delta u-(u\cdot\nabla)u-\nabla p,
\]
\[
\partial_t^2u
=\nu\Delta\partial_tu
-(\partial_tu\cdot\nabla)u
-(u\cdot\nabla)\partial_tu
-\nabla\partial_tp,
\]
\[
\begin{aligned}
\partial_t^3u={}&\nu\Delta\partial_t^2u
-(\partial_t^2u\cdot\nabla)u
-2(\partial_tu\cdot\nabla)\partial_tu \\
&-(u\cdot\nabla)\partial_t^2u
-\nabla\partial_t^2p.
\end{aligned}
\]
The complete pattern is
\[
\partial_t^{\,n+1}u
=\nu\Delta\partial_t^nu
-\sum_{a=0}^{n}\binom na
\bigl(\partial_t^au\cdot\nabla\bigr)\partial_t^{\,n-a}u
-\nabla\partial_t^np.
\]
Only now set
\[
V_n:=\partial_t^nu,
\qquad
Q_n:=\partial_t^np.
\]
Every rise in temporal order brings a Laplacian of the preceding response. Viscosity links one step up the time ladder to two spatial derivatives. The nonlinear products spread across all lower temporal orders, and pressure rises with them, so no row has been separated from the full equation.

At the critical height, that link has an exact energy form:
\[
E_s'=-2\nu E_{s+1}+\mathcal P_s,
\qquad
\mathcal P_s
:=-\bigl\langle\Lambda^su,
\Lambda^s\bigl((u\cdot\nabla)u+\nabla p\bigr)\bigr\rangle.
\]
Repeated differentiation gives
\[
H^{(n)}
=(-2\nu)^nE_{n+\frac12}
+\sum_{j=0}^{n-1}(-2\nu)^j
\partial_t^{\,n-1-j}\mathcal P_{j+\frac12}.
\]
The \(n\)-th temporal response contains the spatial energy at height \(n+\tfrac12\), together with every differentiated nonlinear and pressure transfer. As the temporal order rises, the Sobolev height rises with it.

In three dimensions,
\[
s>k+\frac32
\quad\Longrightarrow\quad
H^s(\mathbb R^3)\hookrightarrow C^k(\mathbb R^3),
\]
and hence
\[
\bigcap_{s<\infty}H^s(\mathbb R^3)
=H^\infty(\mathbb R^3)
\hookrightarrow C^\infty(\mathbb R^3).
\]
A blow-up would demand a loss of spatial regularity as the temporal responses rise. The viscous relation points toward higher spatial regularity instead. The proof succeeds if it can carry that relation through every finite row without breaking compatibility between the rows.

It does this one finite block at a time. For each finite temporal order \(N\) and spatial order \(m\), the whole-terminal mixed estimate supplies
\[
V_N\in L^2(0,T_*;H^{m+1}),
\qquad
V_{N+1}\in L^2(0,T_*;H^{m-1}),
\]
and, in its full finite norm \(\mathcal R_{N,m}\),
\[
\sup_{0<T<T_*}\mathcal R_{N,m}(u;T)<\infty.
\]
The constant may depend on \(N\) and \(m\). No estimate has to remain uniform as both indices rise. Each finite block keeps the nonlinear products, pressure, viscosity, and adjacent temporal rows together.

All of these blocks come from one velocity history generated by \(u_0\). Their overlaps already agree, so they fit into one compatible intersection:
\[
u_0
\longrightarrow
\{\mathcal R_{N,m}[u_0]\}_{N,m<\infty}
\longrightarrow
\mathcal I[u_0]
:=
\bigcap_{N,m<\infty}H_t^N(0,T_*;H_x^m).
\]
The proof has now reached every finite spatial height and every finite temporal row, without asking for an infinite-order constant.

For each finite \(m\), the intersection contains the adjacent pair
\[
u\in L^2(0,T_*;H^{m+1}),
\qquad
u_t\in L^2(0,T_*;H^{m-1}).
\]
The Hilbert-space trace theorem gives
\[
u\in C([0,T_*];H^m).
\]
Because this holds for every \(m\), and every trace belongs to one velocity field,
\[
u_*:=u(T_*)\in\bigcap_{m<\infty}H^m=H^\infty.
\]
The pressure and velocity recurrences then rebuild the complete endpoint jet:
\[
Q_N
=R_iR_j\sum_{a=0}^{N}\binom NaV_{a,i}V_{N-a,j},
\]
\[
V_{N+1}
=\nu\Delta V_N
-\sum_{a=0}^{N}\binom Na(V_a\cdot\nabla)V_{N-a}
-\nabla Q_N.
\]
Every \(V_N\) and \(Q_N\) lies in \(H^\infty\) at \(T_*\). Smooth local existence restarts the solution from \(u_*\), and uniqueness glues that restart to the original velocity history. The solution continues past \(T_*\), contradicting the assumption that \(T_*\) was maximal. Thus the finite terminal time cannot occur.

This is not an a priori estimate proof. It is an intersection-to-estimate proof. Viscosity relates the velocity's time derivatives to its spatial Sobolev derivatives. Taken through every finite order, those two directions meet in a common intersection. At that intersection, \(H^\infty(\mathbb R^3)\hookrightarrow C^\infty(\mathbb R^3)\) gives smoothness, and the familiar estimates can be taken from whichever finite Sobolev height they require. The estimates come from the intersection, not the other way around.

\end{document}
